\documentclass[11pt]{amsart}

\usepackage[margin=1.1in]{geometry}
\usepackage{amsmath,amssymb,amsthm}
\usepackage{enumitem}
\usepackage{graphicx}
\usepackage[colorlinks=true,linkcolor=blue,citecolor=blue,urlcolor=blue]{hyperref}

\newtheorem{theorem}{Theorem}
\newtheorem{lemma}{Lemma}

\newcommand{\lip}{\operatorname{lip}}
\newcommand{\Lip}{\operatorname{Lip}}
\newcommand{\eps}{\varepsilon}
\newcommand{\R}{\mathbb{R}}

\title[Proper purely unrectifiable spaces]{Proper purely 1-unrectifiable metric spaces have locally flat uniform separation}
\author{arxiv\_automatic\_math run fa\_banach\_001}
\date{2026-06-29}

\begin{document}
\maketitle

\section*{Status}

Result type: partial result.  Status: likely valid pending human review.

This packet addresses Question~2.7 (TeX label \texttt{q:spu p1u}) in
Aliaga's exposition \cite{Aliaga2606}.  The question asks whether every purely
$1$-unrectifiable metric space $M$ has $\lip(M)$ separating points uniformly.
The packet proves an affirmative answer for every proper purely
$1$-unrectifiable metric space.  Here proper means that every closed bounded
ball is compact.  This is still not the full arbitrary metric-space case.

\section*{Open-problem evidence}

The source question appears on page 6 of the rendered source PDF:

\begin{center}
\includegraphics[width=\textwidth]{figures/open_problem_crop.png}
\end{center}

Transcription: Let $M$ be a purely $1$-unrectifiable metric space. Does
$\lip(M)$ separate points of $M$ uniformly?

\section*{Definitions}

A metric space is proper if every closed bounded ball is compact.  A
real-valued Lipschitz function $f$ on $M$ is locally flat at $p\in M$ if its
asymptotic Lipschitz constant at $p$ is zero; equivalently, the small-scale
Lipschitz constants of $f$ at $p$ tend to zero.  We write $\lip(M)$ for the
space of Lipschitz functions that are locally flat at every point.  The space
$\lip(M)$ separates points uniformly if for every $a,b\in M$ and every
$\eps>0$ there is $h\in\lip(M)$ with $\Lip(h)\leq1$ and
\[
  h(b)-h(a)>d(a,b)-\eps.
\]

\section*{Input from the compact theory}

We use the compact approximation lemma quoted in \cite[Lemma 2.11]{Aliaga2606}
from Bate's work \cite{Bate}.  In the form needed here, if $K$ is compact and
purely $1$-unrectifiable, $f\in\Lip(K)$, and $\eta>0$, then there are
$g\in\Lip(K)$ and $\delta>0$ such that
\begin{enumerate}[label=(\roman*),leftmargin=*]
\item $|g(z)-f(z)|\leq\eta$ for all $z\in K$;
\item for all $u\neq v$ in $K$,
\[
  \frac{|g(u)-g(v)|}{d(u,v)}
  \leq
  \frac{|f(u)-f(v)|}{d(u,v)}+\eta;
\]
\item for all $u\neq v$ in $K$ with $d(u,v)\leq\delta$,
\[
  \frac{|g(u)-g(v)|}{d(u,v)}\leq\eta.
\]
\end{enumerate}

\section*{Proof intuition}

The compact AGPP theorem is proved by repeatedly applying the approximation
lemma until small-scale slopes are forced to zero.  For a proper noncompact
space, closed balls are compact, so we can run the same idea on larger and
larger balls.  After improving a function on a compact ball, extend it back to
the whole space by McShane.  Later improvements on larger balls may change the
old ball, but the lemma's slope inequality says they can increase old slopes
only by the chosen error.  Choosing summable errors with tails tending to zero
makes the final diagonal limit locally flat at every point, because every
point eventually sits in the interior of some compact ball.

\section*{Main theorem}

\begin{theorem}
\label{thm:proper}
Let $M$ be a proper purely $1$-unrectifiable metric space.  Then $\lip(M)$
separates points of $M$ uniformly.
\end{theorem}

\begin{proof}
Fix $a,b\in M$ and $\eps>0$.  The case $a=b$ is trivial, so put
$D=d(a,b)>0$.  We will construct a locally flat function whose Lipschitz
constant is slightly larger than $1$, then rescale.

Choose a summable sequence $(\eta_n)_{n\geq 1}$ of positive numbers so small
that
\[
  S:=\sum_{n=1}^\infty \eta_n
  < \frac{\eps}{D+3}.
\]
Let $o=a$ and choose radii $R_n\uparrow\infty$ with $b\in K_1$, where
\[
  K_n=\overline{B}(o,R_n).
\]
Each $K_n$ is compact by properness and purely $1$-unrectifiable as a subset
of $M$.

Set $F_0(z)=d(a,z)$ on $M$.  Inductively suppose that $F_{n-1}$ is defined on
$M$.  Apply the compact approximation lemma to
$F_{n-1}|_{K_n}\in\Lip(K_n)$ with parameter $\eta_n$.  This gives
$g_n\in\Lip(K_n)$ and $\delta_n>0$ such that
\begin{align}
  |g_n(z)-F_{n-1}(z)|&\leq\eta_n &&(z\in K_n), \label{eq:uniform-step}\\
  \frac{|g_n(u)-g_n(v)|}{d(u,v)}
  &\leq
  \frac{|F_{n-1}(u)-F_{n-1}(v)|}{d(u,v)}+\eta_n
  &&(u\neq v\in K_n), \label{eq:slope-step}\\
  \frac{|g_n(u)-g_n(v)|}{d(u,v)}
  &\leq\eta_n
  &&(u\neq v\in K_n,\ d(u,v)\leq\delta_n). \label{eq:flat-step}
\end{align}
Let $F_n$ be any McShane extension of $g_n$ to $M$.  Then
$F_n|_{K_n}=g_n$ and
\[
  \Lip(F_n)\leq\Lip(g_n)\leq \Lip(F_{n-1})+\eta_n.
\]
Consequently $\Lip(F_n)\leq 1+\sum_{k=1}^n\eta_k$.

For a fixed $N$, the restrictions $F_n|_{K_N}$ are uniformly Cauchy for
$n\geq N$.  Indeed, if $n>N$, then $K_N\subset K_n$, and
\eqref{eq:uniform-step} gives
\[
  \sup_{z\in K_N}|F_n(z)-F_{n-1}(z)|\leq\eta_n.
\]
Thus the restrictions converge uniformly on every $K_N$ to a compatible
function $F:M\to\R$.  The Lipschitz estimates pass to the limit and give
\[
  \Lip(F)\leq 1+S.
\]
Since $a,b\in K_1$, telescoping \eqref{eq:uniform-step} also gives
\[
  |F(a)-F_0(a)|\leq S,\qquad |F(b)-F_0(b)|\leq S.
\]
Therefore
\[
  F(b)-F(a)\geq D-2S.
\]

We now prove that $F$ is locally flat.  Fix $p\in M$ and $\alpha>0$.  Choose
$N$ so large that $p$ lies in the interior of $K_N$ and
\[
  \sum_{n=N}^\infty \eta_n < \alpha.
\]
Let $\rho>0$ be such that $B(p,\rho)\subset K_N$.  For $u,v\in K_N$ with
$u\neq v$ and $d(u,v)\leq\delta_N$, \eqref{eq:flat-step} gives the slope bound
$\eta_N$ for $F_N$.  For every $m>N$, applying \eqref{eq:slope-step} on
$K_m$ to the same pair gives
\[
  \frac{|F_m(u)-F_m(v)|}{d(u,v)}
  \leq
  \frac{|F_{m-1}(u)-F_{m-1}(v)|}{d(u,v)}+\eta_m.
\]
Letting $m\to\infty$ yields
\[
  \frac{|F(u)-F(v)|}{d(u,v)}
  \leq \sum_{n=N}^\infty \eta_n <\alpha
\]
whenever $u,v\in K_N$ and $d(u,v)\leq\delta_N$.  In particular this holds for
$u,v\in B(p,\rho/2)$ with $d(u,v)\leq\delta_N$.  Hence the asymptotic
Lipschitz constant of $F$ at $p$ is $0$.  Since $p$ was arbitrary,
$F\in\lip(M)$.

Finally set $H=F/(1+S)$.  Then $H\in\lip(M)$ and $\Lip(H)\leq1$.  Moreover
\[
  H(b)-H(a)
  \geq \frac{D-2S}{1+S}
  = D-\frac{S(D+2)}{1+S}
  > D-\eps,
\]
by the choice of $S$.  Thus $\lip(M)$ separates the pair $a,b$ uniformly, and
the pair was arbitrary.
\end{proof}

\section*{Verification notes}

The proof is non-computational.  The main things to check are:
\begin{itemize}[leftmargin=*]
\item The compact approximation lemma is applied only to compact balls, which
  are purely $1$-unrectifiable because pure unrectifiability passes to
  subsets.
\item McShane extension is used only after the compact modification; all later
  estimates on older compact balls are made before extension matters, because
  the next compact ball contains the old one.
\item Local flatness follows from the summable tail of the slope errors and
  from the fact that every point eventually lies in the interior of a compact
  ball in the exhaustion.
\end{itemize}

\section*{Novelty and limitations}

The run indexes were searched for \texttt{2606.02918}, \texttt{q:spu},
\texttt{locally flat}, \texttt{separates points uniformly}, \texttt{proper},
\texttt{boundedly compact}, and \texttt{purely 1-unrectifiable}.  The only
same-question packet found was the earlier metric-wedge partial result from
this run.  A bounded web phrase search on 2026-06-29 for proper/boundedly
compact variants of the locally-flat separation question found no exact
proper-space closure result.

This is not a full answer to Question~2.7.  The proof uses compact
neighborhoods essentially: without local compactness, the exhaustion does not
place small neighborhoods of a point inside a compact pure set, and the
argument gives no local-flatness control.

\section*{Human review recommendation}

Reviewers should focus on the diagonal/exhaustion step: the restrictions
$F_n|_{K_N}$ are Cauchy, and the slope estimates on $K_N$ survive all later
steps by the one-sided slope inequality in Bate's compact lemma.  If that
bookkeeping is accepted, the theorem is a direct noncompact extension of the
compact approximation proof.

\begin{thebibliography}{9}

\bibitem{Aliaga2606}
R. J. Aliaga,
\emph{Lipschitz-free spaces and purely 1-unrectifiable metric spaces},
arXiv:2606.02918, 2026.
\url{https://arxiv.org/abs/2606.02918}

\bibitem{Bate}
D. Bate,
\emph{Purely unrectifiable metric spaces and perturbations of Lipschitz functions},
Acta Math. \textbf{224} (2020), 1--65.

\bibitem{AGPP}
R. J. Aliaga, C. Gartland, C. Petitjean, and A. Prochazka,
\emph{Purely 1-unrectifiable metric spaces and locally flat Lipschitz functions},
Trans. Amer. Math. Soc. \textbf{375} (2022), 3529--3567.
\url{https://arxiv.org/abs/2103.09370}

\end{thebibliography}

\end{document}
